\begin{table}[!htbp]
\centering
\caption{Estimated non-inferiority power in the SWIFT DIRECT-inspired illustration.}
\label{tab:swift-direct-siga-summary}
\begin{threeparttable}
\small
\setlength{\tabcolsep}{5.5pt}
\renewcommand{\arraystretch}{1.10}
\begin{tabular}{lrrrrr}
\toprule
Score & SIGA-S (\%) & SIGA-R (\%) & RT (\%) & R$-$RT (pp) & Mean $\widehat\rho_n$ \\ 
\midrule
Unadjusted & 81.48 & 81.44 & 81.28 & +0.16 & 1.003 \\
Baseline-adjusted & 81.53 & 81.48 & 81.50 & -0.02 & 1.003 \\
\bottomrule
\end{tabular}
\begin{tablenotes}[flushleft]
\footnotesize
\item RT denotes the reference fixed-score randomization test. R$-$RT is the paired difference between SIGA-R and RT rejection probabilities in percentage points. The diagnostic $\widehat\rho_n$ is the raw SIGA-R variance divided by the SIGA-S variance.
\item All three procedures were evaluated on the same 100,000 outer trials. SIGA-S targets repeated-sampling power for the marginal risk difference, whereas SIGA-R targets the rejection probability of the prespecified reference randomization test.
\end{tablenotes}
\end{threeparttable}
\end{table}

